% Source: physical PDF pp. 419--430 (printed pp. 396--407). Coverage:
% complete Section 22.6, from its heading through immediately before
% Section 22.7; Eqs. (22.6.1)--(22.6.35), no figures or tables.
\subsection{Consistency Conditions}
\label{sec:22.6}

The numerical coefficient appearing in the anomaly for any symmetry
depends on the matter content of the theory. On the other hand, the
\emph{form} of the anomaly is largely independent of the details of the
theory, because it is governed by consistency conditions, first presented
in 1971 by Wess and Zumino~\hyperref[ch22-ref-18]{[18]}.

Even when we are interested in anomalies in global symmetry currents, in
order to derive the consistency conditions it is convenient to imagine
that all symmetry currents are coupled to gauge fields, which for
non-Abelian symmetries also couple to each other, in such a way that these
symmetries become \emph{local} symmetries of the Lagrangian density. We can
always return to the case of global symmetry by letting the corresponding
gauge coupling constants become infinitesimal. Apart from anomalies, the
effective action $\Gamma[A]$ in a background gauge field
$A_{a\mu}(x)$ will in this formalism be invariant under infinitesimal
transformations
$A_{b\mu}(y)\rightarrow A_{b\mu}(y)
+i\int d^4x\,\epsilon_a(x)\mathcal T_a(x)A_{b\mu}(y)$
on the gauge field, where in order to reproduce the transformation
\eqref{eq:15.1.9} we must take
\begin{equation}
 -i\mathcal T_a(x)
 =-\frac{\partial}{\partial x^\mu}
 \frac{\delta}{\delta A_{a\mu}(x)}
 -C_{abc}A_{b\mu}(x)
 \frac{\delta}{\delta A_{c\mu}(x)}.
 \tag{22.6.1}\label{eq:22.6.1}
\end{equation}
Taking anomalies into account, $\mathcal T_a(x)$ no longer annihilates
$\Gamma[A]$, but rather
\begin{equation}
 \mathcal T_a(x)\Gamma[A]=G_a[x;A],
 \tag{22.6.2}\label{eq:22.6.2}
\end{equation}
where $G_a[x;A]$ represents the effect of the anomaly.
Eq.~\eqref{eq:22.6.2} may also be written as a formula for the covariant
divergence of the expectation value of the current:
\begin{equation}
 D_\mu\langle J^\mu_a(x)\rangle=-iG_a[x;A],
 \tag{22.6.3}\label{eq:22.6.3}
\end{equation}
where
\begin{equation}
 \langle J^\mu_a(x)\rangle
 =\frac{\delta}{\delta A_{a\mu}(x)}\Gamma[A]
 \tag{22.6.4}\label{eq:22.6.4}
\end{equation}
and $D_\mu$ is the gauge-covariant derivative \eqref{eq:15.1.10}, taken
here in the adjoint representation with
$(t_b)_{ca}=-iC_{abc}$.

The Wess--Zumino consistency conditions follow from the commutation
relations\footnote{The factor $-i$ was inserted in
Eq.~\eqref{eq:22.6.1} in order to provide the conventional factor $+i$
accompanying the structure constant $C_{abc}$ in this
commutation relation. A reminder: we are using a basis for the Lie algebra
in which the structure constants are totally antisymmetric, and so we do
not distinguish between upper and lower gauge indices.}
\begin{equation}
 [\mathcal T_a(x),\mathcal T_b(y)]
 =iC_{abc}\delta^4(x-y)\mathcal T_c(x).
 \tag{22.6.5}\label{eq:22.6.5}
\end{equation}
From Eqs.~\eqref{eq:22.6.2} and \eqref{eq:22.6.3}, we derive the general
consistency condition
\begin{equation}
 \mathcal T_a(x)G_b[y;A]
 -\mathcal T_b(x)G_a[y;A]
 =iC_{abc}\delta^4(x-y)G_c[y;A].
 \tag{22.6.6}\label{eq:22.6.6}
\end{equation}

These consistency conditions were originally derived by Wess and Zumino
for the chiral $SU(3)\times SU(3)$ symmetry of the strong interactions, a
special case of physical as well as historical importance. Here the
generators $\mathcal T_a(x)$ acting on the gauge fields consist of the
even parity generators $\mathcal V_a(x)$ defined by
Eq.~\eqref{eq:22.3.31}, and the odd parity generators $\mathcal A_a(x)$
defined by \eqref{eq:22.3.32}.\footnote{Another reminder: as mentioned in
\hyperref[sec:22.3]{Section 22.3}, the $\mathcal A_a$ and $\mathcal V_a$
used here are what were called $Y_a$ and $X_a$ in
Ref.~\hyperref[ch22-ref-8]{[8]}.} They satisfy the commutation relations
\[
 [\mathcal V_a(x),\mathcal V_b(y)]
 =i\delta^4(x-y)f_{abc}\mathcal V_c(x),
\]
\[
 [\mathcal V_a(x),\mathcal A_b(y)]
 =i\delta^4(x-y)f_{abc}\mathcal A_c(x),
\]
\[
 [\mathcal A_a(x),\mathcal A_b(y)]
 =i\delta^4(x-y)f_{abc}\mathcal V_c(x),
\]
where $f_{abc}$ are the $SU(3)$ structure constants. Since the $SU(3)$
subgroup generated by the $\mathcal V_a$ is not spontaneously broken, it
is convenient to treat the integration over fermion momenta in such a way
as to preserve invariance under gauge transformations generated by
$\mathcal V_a$, so that
\[
 \mathcal V_a(x)\Gamma=0,
\]
leaving us with the non-zero anomaly
\[
 \mathcal A_a(x)\Gamma=G_a(x).
\]
The non-trivial consistency conditions are then
\[
 \mathcal V_a(x)G_b(y)
 =i\delta^4(x-y)f_{abc}G_c(x)
\]
and
\[
 \mathcal A_a(x)G_b(y)-\mathcal A_b(x)G_a(y)=0.
\]
The first of these simply says that $G_a(x)$ transforms like an octet
under ordinary $SU(3)$ transformations. The second condition imposes
other strong constraints on $G_a(x)$. The reader may check that this
condition is satisfied by the Bardeen formula \eqref{eq:22.3.34} for
$G_a(x)$. We will not go into this here, but will instead use as an
illustration a general gauge theory with all currents treated symmetrically.
In \hyperref[sec:22.3]{Section 22.3} we quoted the formula
\eqref{eq:22.3.38} for the anomaly in this case, but did not derive the
terms in this formula of higher than second order in the gauge fields. Here
we shall show that these terms are dictated by the consistency conditions
\eqref{eq:22.6.6}.

For this purpose, and also to allow for further generalizations, it is very
convenient to reformulate the set of Wess--Zumino consistency conditions
as a condition of invariance under the BRST transformations described in
\hyperref[sec:15.7]{Section 15.7}. Let us introduce a ghost field
$\omega_a$, and define the nilpotent BRST operator $s$ for a general
gauge theory by
\begin{equation}
 sA_{a\mu}
 =\partial_\mu\omega_a
 +C_{abc}A_{b\mu}\omega_c,
 \tag{22.6.7}\label{eq:22.6.7}
\end{equation}
\begin{equation}
 s\omega_a
 =-\frac12 C_{abc}\omega_b\omega_c,
 \tag{22.6.8}\label{eq:22.6.8}
\end{equation}
it being understood that $s$ satisfies the distributive rule
$s(AB)=(sA)B\pm A(sB)$, the sign being negative where $A$ is a fermionic
quantity like $\omega_a$, and otherwise positive. In place of the
anomaly function $G_a[x;A]$, we shall work with a functional
\begin{equation}
 G[\omega,A]
 =\int\omega_a(x)G_a[x;A]\,d^4x.
 \tag{22.6.9}\label{eq:22.6.9}
\end{equation}
Then (keeping in mind that $\omega_a$ is fermionic)
\begin{align*}
 sG[\omega,A]
 ={}&-\frac12 C_{abc}
 \int d^4x\,\omega_b(x)\omega_c(x)G_a[x;A]
 \\
 ={}&\int d^4x\,\omega_a(x)\int d^4y
 \left[
 \frac{\partial\omega_b(y)}{\partial y^\mu}
 +C_{bcd}A_{c\mu}(y)\omega_d(y)
 \right]
 \frac{\delta G_a[x;A]}{\delta A_{b\mu}(y)}
 \\
 ={}&\int d^4x\int d^4y\,\omega_a(x)\omega_b(y)
 \left[
 -\frac12 C_{abc}\delta^4(x-y)G_c[x;A]
 +i\mathcal T_b(y)G_a[x;A]
 \right].
\end{align*}
Because the ghost fields anticommute with each other, this can be written
\[
 \begin{aligned}
 sG[\omega,A]
 =-\frac12 i\int d^4x\int d^4y\,
 \omega_a(x)\omega_b(y)
 \Bigl[
 &iC_{abc}\delta^4(x-y)G_c[x;A]
 \\
 &+\mathcal T_b(y)G_a[x;A]
 -\mathcal T_a(x)G_b(y)
 \Bigr].
 \end{aligned}
\]
We see that the consistency condition \eqref{eq:22.6.6} will hold if and
only if $G[\omega,A]$ is BRST-invariant
\begin{equation}
 sG[\omega,A]=0
 \tag{22.6.10}\label{eq:22.6.10}
\end{equation}
for all ghost fields $\omega_a(x)$.

Now consider the possibility that the anomaly $G[A;\omega]$ could be
written as the BRST operator $s$ acting on a local functional $F[A]$:
\begin{equation}
 G[A;\omega]=sF[A].
 \tag{22.6.11}\label{eq:22.6.11}
\end{equation}
(Note that the functional $F$ is necessarily independent of the ghost
field, because the operator $s$ adds one ghost field factor, and the
anomaly functional $G$ is already linear in the ghost field.) The BRST
operator satisfies $s^2=0$, so such an anomaly would satisfy the
consistency condition $sG=0$. If $F[A]$ is a \emph{local}
functional\footnote{By a local functional is meant the integral of a local
function, that is, a function of fields and field derivatives at a given
point.} of the gauge field, then it could be subtracted from the action,
thus cancelling the anomaly. The same is true of any term in the anomaly
that can be written as the BRST operator $s$ acting on a local functional;
such a term satisfies the consistency condition by itself, and can be
cancelled by adding a local term to the action. The possible anomalies
that interest us are thus the local functionals $G[\omega;A]$ of ghost
number unity that satisfy the consistency condition \eqref{eq:22.6.10},
modulo terms that can be expressed as $s$ acting on some local functional
of ghost number zero. In accordance with the usual terminology for
nilpotent operators, the equivalence classes of such functionals form what
is called the \emph{cohomology} of the $s$ operator at ghost number unity.

We can also express this in terms of the local densities themselves. We
can write the anomaly (or any term in the anomaly) as
$G=\int d^4x\,\mathcal G(x)$, where $\mathcal G(x)$ is a power series in
the gauge and ghost fields and their derivatives at the spacetime point
$x$. The condition $sG=0$ then is equivalent to the statement that
\begin{equation}
 s\mathcal G(x)=\partial_\mu\mathcal F^\mu(x)
 \tag{22.6.12}\label{eq:22.6.12}
\end{equation}
for some function $\mathcal F^\mu(x)$ of fields and field derivatives.
Likewise, the terms in $\mathcal G$ that can be cancelled by adding local
terms to the action are those that are of the form $s\mathcal F$ up to
possible derivatives. Thus the anomalies that interest us are the local
functions of ghost number unity that satisfy the consistency condition
\eqref{eq:22.6.12}, modulo terms that can be expressed as $s$ acting on
some local functional of ghost number zero, modulo possible derivatives.
This is known as the cohomology of $s$ at ghost number unity in the space
of local functions, modulo derivatives, and denoted $H^1(s\mathbin{|}d)$.

Algebraic methods have been used to work out the cohomology of the BRST
operator $s$, and thereby deduce the form of the anomaly in general gauge
theories~\hyperref[ch22-ref-19]{[19]}. This approach leaves unknown only
constant coefficients that depend on the matter content of the theory and
need to be calculated by the methods of Sections 22.2 or 22.3. Since we
have already calculated the terms in the anomaly of second order in the
gauge fields for general gauge theories, including their constant
coefficients, here we will use the
consistency condition \eqref{eq:22.6.12} to calculate the terms of higher
order in the fields.

We saw in \hyperref[sec:22.3]{Section 22.3} that when all currents are
treated symmetrically, the terms of second order in the gauge fields are
one-third the expression \eqref{eq:22.3.26}. This is an operator of
dimensionality four (in units of mass), while the Wess--Zumino consistency
condition \eqref{eq:22.6.6} relates only operators of the same
dimensionality, so to satisfy this condition we should add only terms of
higher order in the gauge field that have the same dimensionality. We
therefore seek a solution of the consistency conditions in the (not
necessarily unique) form
\begin{equation}
 \begin{aligned}
 G_a
 \equiv i[\langle\partial_\mu J^\mu_a\rangle]_{\rm anom}
 =-\frac{i}{24\pi^2}\epsilon^{\kappa\nu\lambda\rho}
 \operatorname{Tr}\Bigl\{T_a\bigl[
 &\partial_\kappa A_\nu\,\partial_\lambda A_\rho
 +ic_1\partial_\kappa A_\nu A_\lambda A_\rho
 \\
 &+ic_2A_\kappa\partial_\nu A_\lambda A_\rho
 +ic_3A_\kappa A_\nu\partial_\lambda A_\rho
 \\
 &-c_4A_\kappa A_\nu A_\lambda A_\rho
 \bigr]\Bigr\},
 \end{aligned}
 \tag{22.6.13}\label{eq:22.6.13}
\end{equation}
where $A_\mu\equiv A_{a\mu}T_a$, and the $c_i$ are constants to
be determined.

In order to save a great deal of effort, it will be convenient to rewrite
this in the language of differential forms. (See
\hyperref[sec:8.8]{Section 8.8}.) We introduce a set of c-number parameters
$dx^\mu$ that are taken to anticommute with themselves and all fermionic
fields, such as the ghost field $\omega_a$. The $dx^\mu$ then also
anticommute with the BRST operator $s$. Because
$dx^\kappa dx^\nu dx^\lambda dx^\rho$ is totally antisymmetric, it may be
written as
\begin{equation}
 dx^\kappa dx^\nu dx^\lambda dx^\rho
 =\epsilon^{\kappa\nu\lambda\rho}d^4x,
 \qquad
 d^4x=dx^0dx^1dx^2dx^3.
 \tag{22.6.14}\label{eq:22.6.14}
\end{equation}
We also introduce the exterior derivative
\[
 d\equiv dx^\mu\frac{\partial}{\partial x^\mu},
\]
which since derivatives commute is nilpotent as well as anticommuting
with $s$:
\begin{equation}
 d^2=0,
 \qquad
 ds+sd=0.
 \tag{22.6.15}\label{eq:22.6.15}
\end{equation}
Finally, we introduce the anticommuting quantities
\begin{equation}
 A\equiv iA_\mu dx^\mu
 =iA_{a\mu}T_a dx^\mu,
 \qquad
 \omega\equiv i\omega_a T_a.
 \tag{22.6.16}\label{eq:22.6.16}
\end{equation}
In this notation, Eq.~\eqref{eq:22.6.13} reads
\begin{equation}
 \begin{aligned}
 G[\omega,A]
 =\frac{1}{24\pi^2}\int\operatorname{Tr}\Bigl\{
 \omega\bigl[
 &(dA)^2+c_1(dA)A^2+c_2A(dA)A
 \\
 &+c_3A^2(dA)+c_4A^4
 \bigr]\Bigr\}.
 \end{aligned}
 \tag{22.6.17}\label{eq:22.6.17}
\end{equation}
In order to implement the consistency condition \eqref{eq:22.6.10}, we
note that the BRST transformation rules \eqref{eq:22.6.7} and
\eqref{eq:22.6.8} may be written as
\begin{equation}
 sA=-d\omega+\{A,\omega\},
 \tag{22.6.18}\label{eq:22.6.18}
\end{equation}
\begin{equation}
 s\omega=\omega^2.
 \tag{22.6.19}\label{eq:22.6.19}
\end{equation}
Now, the BRST transformation of the last term in Eq.~\eqref{eq:22.6.17}
is given by
\[
 \begin{aligned}
 s\operatorname{Tr}[\omega A^4]
 =\operatorname{Tr}\bigl[
 &\omega^2A^4-\omega\{A,\omega\}A^3
 +\omega A\{A,\omega\}A^2
 \\
 &-\omega A^2\{A,\omega\}A
 +\omega A^3\{A,\omega\}
 \bigr]
 +\text{$\omega d\omega A^3$ terms}
 \\
 ={}&-\operatorname{Tr}[\omega^2A^4]
 +\text{$\omega d\omega A^3$ terms}.
 \end{aligned}
\]
There is no other contribution to $sG$ proportional to
$\operatorname{Tr}[\omega^2A^4]$, so the consistency condition
\eqref{eq:22.6.10} can only be satisfied if $c_4=0$. With $c_4=0$, a
straightforward calculation gives
\[
 \begin{aligned}
 sG=\frac{1}{24\pi^2}\int\operatorname{Tr}\Bigl\{
 &-(dA)^2\omega^2+\omega d\omega A dA-d\omega\,\omega dA A
 \\
 &-A\omega dA d\omega-\omega A d\omega dA
 \\
 &+c_1[\omega dA d\omega A-\omega dA A d\omega]
 \\
 &+c_2[\omega d\omega dA A-\omega A dA d\omega]
 \\
 &+c_3[\omega d\omega A dA-\omega A d\omega dA]
 \\
 &-c_1[-\omega A d\omega A^2+\omega d\omega A^3
       +\omega dA A^2\omega]
 \\
 &-c_2[\omega A^2d\omega A-\omega A d\omega A^2
       +\omega A dA A\omega]
 \\
 &-c_3[-\omega A^3d\omega+\omega A^2d\omega A
       +\omega A^2dA\omega]\Bigr\}.
 \end{aligned}
\]
We do not need to assume that the integrand vanishes, but only that it is
the derivative of some local function, so that its integral vanishes. This
condition must be satisfied separately for the terms involving two
derivatives and those involving one derivative, since no cancellation can
occur between terms with different numbers of derivatives. It is not hard
to see that the terms in the integrand involving just one derivative are
of the form $d\mathcal F$ if we take
$c_1=-c_2=+c_3=c$. The remaining terms are a total derivative if
$c=-1/2$, thus justifying the previously quoted result
\eqref{eq:22.3.38}. This result is often expressed more compactly as
\begin{equation}
 \begin{aligned}
 G[\omega,A]
 &=\frac{1}{24\pi^2}\int
 \operatorname{Tr}\left\{\omega\,d\left[
 A\,dA-\frac12A^3\right]\right\}
 \\
 &=\frac{1}{24\pi^2}\int
 \operatorname{Tr}\left\{\omega\,d\left[
 AF+\frac12A^3\right]\right\},
 \end{aligned}
 \tag{22.6.20}\label{eq:22.6.20}
\end{equation}
where $F$ is the matrix-valued field-strength two-form
\begin{equation}
 F\equiv\frac12 it_a F_{a\mu\nu}\,dx^\mu dx^\nu
 =dA-A^2.
 \tag{22.6.21}\label{eq:22.6.21}
\end{equation}
The anomaly does not have to be put in the form \eqref{eq:22.6.13}, so
Eq.~\eqref{eq:22.6.20} is not the unique result for $G[\omega,A]$. Results
quoted in the next section show that, for any subgroup $H$ of $G$ for
which $\operatorname{Tr}\{t_i\{t_j,t_k\}\}=0$ for all generators of $H$,
it is possible to add local terms to the action in such a way that the
anomaly $G_i$ vanishes when $t_i$ is any generator of $H$.

There is an elegant algebraic tool, known as the
\emph{Stora--Zumino descent equations},\hyperref[ch22-ref-18a]{[18a]} for
constructing a solution of the consistency conditions. It is just as easy
to describe this method in a spacetime of any even dimensionality as in
four spacetime dimensions, so we will take the spacetime dimensionality to
be $2n$. To start, one must imagine that at least two additional variables
have been added to the $2n$ coordinates of space and time, in order to give
meaning to the $(2n+2)$-form $\operatorname{Tr}F^{n+1}$. Note that
\begin{equation}
 dF=-d(A^2)=-(dA)A+A(dA)=[A,F],
 \tag{22.6.22}\label{eq:22.6.22}
\end{equation}
so that $\operatorname{Tr}F^{n+1}$ is closed:
\begin{equation}
 d\operatorname{Tr}F^{n+1}
 =(n+1)\operatorname{Tr}\{(dF)F^n\}
 =\operatorname{Tr}\{[A,F^{n+1}]\}=0.
 \tag{22.6.23}\label{eq:22.6.23}
\end{equation}
As long as the extended spacetime is simply connected, Poincaré's theorem
then tells us that $\operatorname{Tr}\{F^{n+1}\}$ is exact, in the sense
that there is a $(2n+1)$-form $\Omega_{2n+1}$ (known as the
\emph{Chern--Simons form}), for which
\begin{equation}
 \operatorname{Tr}\{F^{2n+2}\}=d\Omega_{2n+1}.
 \tag{22.6.24}\label{eq:22.6.24}
\end{equation}
Further, $\operatorname{Tr}\{F^{n+1}\}$ is manifestly gauge-invariant and
depends only on the gauge field, so it is BRST-invariant:
\begin{equation}
 s\operatorname{Tr}\{F^{n+1}\}=0.
 \tag{22.6.25}\label{eq:22.6.25}
\end{equation}
The `differentials' $dx^\mu$ are understood to anticommute with fermionic
fields like the ghost field $\omega_a$, so the operator $d$
anticommutes with the operator $s$ defined by Eqs.~\eqref{eq:22.6.7} and
\eqref{eq:22.6.8}:
\[
 sd+ds=0.
\]
Because $s$ is nilpotent, it follows then that $s\Omega_{2n+1}$ is also
closed:
\[
 d(s\Omega_{2n+1})=-s\operatorname{Tr}\{F^{n+1}\}=0.
\]
Again using Poincaré's theorem, this means that there must be a $2n$-form
$\Omega^1_{2n}$, of first order in $\omega_a$, for which
\begin{equation}
 s\Omega_{2n+1}=d\Omega^1_{2n}.
 \tag{22.6.26}\label{eq:22.6.26}
\end{equation}
Furthermore, $d(s\Omega^1_{2n})=-s^2\Omega_{2n+1}=0$, so there is also a
$(2n-1)$-form $\Omega^2_{2n-1}$, of second order in the ghost field, for
which
\begin{equation}
 s\Omega^1_{2n}=d\Omega^2_{2n-1}.
 \tag{22.6.27}\label{eq:22.6.27}
\end{equation}
It follows than that the integral of $\Omega^1_{2n}$ over the $2n$
dimensions of space and time is BRST-invariant:
\begin{equation}
 s\int_{\rm spacetime}\Omega^1_{2n}=0,
 \tag{22.6.28}\label{eq:22.6.28}
\end{equation}
even though $\Omega^1_{2n}$ is not itself BRST-invariant. We can thus find
a candidate $\int\Omega^1_{2n}$ for the anomaly functional $G[\omega,A]$
by integrating the two first-order differential equations
$d\Omega_{2n+1}=\operatorname{Tr}\{F^3\}$ and
$d\Omega^1_{2n}=s\Omega_{2n+1}$. General (not unique) solutions of these
equations are
\begin{equation}
 \Omega_{2n+1}
 =(n+1)\int_0^1dt\,\operatorname{Tr}\{AF_t^n\},
 \tag{22.6.29}\label{eq:22.6.29}
\end{equation}
\begin{equation}
 \Omega^1_{2n}
 =n(n+1)\int_0^1dt\,(1-t)
 \operatorname{Tr}\{\omega\,d(AF_t^{n-1})\},
 \tag{22.6.30}\label{eq:22.6.30}
\end{equation}
where $F_t=tF+(t-t^2)A^2$. Evaluation of the integral \eqref{eq:22.6.30}
shows that Eq.~\eqref{eq:22.6.20} gives a result for $G[\omega,A]$
proportional to $\int\Omega^1_4$ in the case of four spacetime dimensions.

We can continue this descent, and derive other useful results. In
particular, it follows from Eq.~\eqref{eq:22.6.27} and the nilpotence of
$s$ that $d(s\Omega^2_{2n-1})=0$, so Poincaré's theorem tells us that
$s\Omega^2_{2n-1}$ is of the form $d\Omega^3_{2n-2}$, so the integral of
$\Omega^2_{2n-1}$ over the $2n-1$ coordinates of space is BRST-invariant:
\begin{equation}
 s\int_{\rm space}\Omega^2_{2n-1}=0.
 \tag{22.6.31}\label{eq:22.6.31}
\end{equation}
Such BRST-invariant functionals of second order in the ghost fields are
candidates~\hyperref[ch22-ref-18b]{[18b]} for so-called
\emph{Schwinger terms}~\hyperref[ch22-ref-18c]{[18c]}. Schwinger terms of
the sort that concern us here arise as anomalous terms
$S_{ab}(\mathbf x,\mathbf y)$ in the equal-time commutation
relations of the time components of two symmetry currents:
\begin{equation}
 [J^0_a(\mathbf x,t),J^0_b(\mathbf y,t)]
 =iC_{abc}J^0_c(\mathbf x,t)
 \delta^{2n-1}(\mathbf x-\mathbf y)
 +S_{ab}(\mathbf x,\mathbf y,t).
 \tag{22.6.32}\label{eq:22.6.32}
\end{equation}
(All operators in this paragraph are taken at the same time $t$, which
will henceforth not be shown explicitly.) From the antisymmetry of the
commutator we have
$S_{ab}(\mathbf x,\mathbf y)
=-S_{ba}(\mathbf y,\mathbf x)$, so all information about
$S_{ab}(\mathbf x,\mathbf y)$ is contained in the functional
\begin{equation}
 S[\omega]\equiv\int d^{2n-1}\mathbf x\,d^{2n-1}\mathbf y\,
 \omega_a(\mathbf x)\omega_b(\mathbf y)
 S_{ab}(\mathbf x,\mathbf y).
 \tag{22.6.33}\label{eq:22.6.33}
\end{equation}
Note that $S_{ab}(\mathbf x,\mathbf y)$ depends in general on
various matter and gauge fields, so $S[\omega]$ generally depends on these
fields as well as on the ghost field $\omega_a(\mathbf x)$. Taking the
commutator of Eq.~\eqref{eq:22.6.32} with a third current
$J^0_c(\mathbf z)$, contracting with
$\omega_a(\mathbf x)\omega_b(\mathbf y)\omega_a(\mathbf z)$,
integrating over $\mathbf x$, $\mathbf y$, and $\mathbf z$, and using the
Jacobi identities, we find
\[
 \begin{aligned}
 0={}&\int d^{2n-1}\mathbf x\,d^{2n-1}\mathbf y\,
 d^{2n-1}\mathbf z\,
 \omega_a(\mathbf x)\omega_b(\mathbf y)
 \omega_a(\mathbf z)
 \\
 &\quad{}\cdot\Bigl[
 iC_{abd}
 S_{cd}(\mathbf z,\mathbf x)
 \delta^{2n-1}(\mathbf x-\mathbf y)
 +[J^0_c(\mathbf z),
 S_{ab}(\mathbf x,\mathbf y)]
 \Bigr].
 \end{aligned}
\]
On functionals of gauge and matter fields like
$S_{ab}(\mathbf x,\mathbf y)$, the action of the BRST operator $s$
is the same as that of a gauge transformation with transformation
parameter $\omega_a$, so
\[
 sS_{ab}(\mathbf x,\mathbf y)
 =i\left[
 \int d^{2n-1}\mathbf z\,\omega_c(\mathbf z)J^0_c(\mathbf z),
 S_{ab}(\mathbf x,\mathbf y)
 \right].
\]
Recalling Eq.~\eqref{eq:22.6.8}, we find that the functional
\eqref{eq:22.6.33} is BRST-invariant
\begin{equation}
 sS[\omega]=0.
 \tag{22.6.34}\label{eq:22.6.34}
\end{equation}
Also, by adding terms to the currents we can change $S[\omega]$ by terms
of the form $sT[\omega]$, so the set of possible Schwinger terms that
could \emph{not} be removed by adding terms to the currents is given by
the cohomology of the BRST operator $s$ at ghost number two---that is, by
BRST-invariant functionals $S$ of second order in the ghost field that are
not themselves of the form $sT$. Eq.~\eqref{eq:22.6.31} shows that
$\int\Omega^2_{2n-1}$ is a candidate for such a functional.

\begin{center}
$* \quad * \quad *$
\end{center}

The analysis of anomalies given so far in this section is strictly
applicable only to anomalies in one-loop order. It is true that a theory
with one-loop anomalies in currents to which quantum gauge fields are
coupled is inconsistent, and therefore does not need to be studied in
higher orders. But the converse does not hold; if a theory with quantum
gauge fields is anomaly-free in one-loop order, we still need to show that
anomalies are absent in higher orders. Also, there is nothing inconsistent
in theories with anomalies in global symmetries, such as the chiral
symmetries of quantum chromodynamics, and for these we need to know
whether higher-order corrections affect the anomalies.

Since BRST transformations act non-linearly on the fields, even in the
absence of anomalies we would not necessarily expect the quantum
effective action $\Gamma[\omega,A]$ to be BRST-invariant beyond the
one-loop approximation. As we saw in
\hyperref[sec:17.1]{Section 17.1}, beyond this approximation we need to
consider functionals not only of gauge and ghost fields but also of their
antifields. (The introduction of antifields is sometimes important even in
one-loop order for another reason: a local functional of fields alone that
satisfies the Wess--Zumino consistency conditions, and that is not
expressible as the BRST operator acting on a local functional of fields
alone, will not be a candidate anomaly if it \emph{can} be expressed as the
antibracket of the action with a local functional of fields and antifields,
because in this case the anomaly can be cancelled by subtracting this term
from the action. This corresponds to a change in the action of the fields
combined with a change in the gauge symmetry obeyed by this action.)

The study of anomalies with antifields included in the action turns out
also to have a cohomological formulation~\hyperref[ch22-ref-20]{[20]},
\hyperref[ch22-ref-21]{[21]}. The analysis of this problem is based on the
Zinn--Justin version \eqref{eq:17.1.10} of what Batalin and Vilkovisky
called the master equation. In the absence of anomalies
$(\Gamma,\Gamma)=0$, so that in one-loop order $(S,\Gamma_1)=0$, where
$S$ is the zeroth-order action and $\Gamma_1$ is the one-loop contribution
to the quantum effective action. In the presence of anomalies we have
instead
\begin{equation}
 (S,\Gamma_1)=G_1,
 \tag{22.6.35}\label{eq:22.6.35}
\end{equation}
where $G_1$ is some functional of fields and antifields which, since $S$
and $\Gamma_1$ have ghost number zero, must have ghost number unity. The
action is assumed to satisfy the classical master equation $(S,S)=0$, so
the antibracket operation in Eq.~\eqref{eq:22.6.35} is nilpotent, and
therefore $(S,G_1)=0$. But if $G_1=(S,F_1)$ for some local functional
$F_1$ of ghost number zero then we can cancel the anomaly to one-loop
order by subtracting the term $F_1$ (which is treated as a quantum
correction of order $\hbar$) from the action and hence from $\Gamma_1$.
(Of course, $G_1=(S,\Gamma_1)$, but in the presence of massless particles
$\Gamma_1$ is not a local functional.) Thus the candidate anomalies are
those local functionals $G_1$ of ghost number unity that are closed, in
the sense that they satisfy $(S,G_1)=0$, but that are not exact, that is,
that cannot be expressed as $G_1=(S,F_1)$ for local functionals $F_1$ of
ghost number unity. In other words, the candidate anomalies correspond to
the cohomology of the antibracket operation $X\mapsto(S,X)$ at ghost
number unity on the space of local functionals of the fields and
antifields.

This is just like the result found earlier in this section, but with the
antibracket $(S,\mathinner{\cdot\cdot\cdot})$ replacing the BRST operator
$s$. If we set the antifields equal to zero in Eq.~\eqref{eq:22.6.35} and
recall that
$(\delta\Gamma_1/\delta\chi_n^*)_{\chi^*=0}=s\chi^n$, we find that the
condition \eqref{eq:22.6.35} in fact yields the condition
$s\Gamma_1=0$, which as we have seen is equivalent to the requirement that
$\Gamma_1$ satisfies the Wess--Zumino consistency conditions. But there
is a sense in which the analysis based on Eq.~\eqref{eq:22.6.35} can be
extended to higher orders.

To see this, suppose that it is found that the antibracket operation
$X\mapsto(S,X)$ has an empty cohomology at ghost number unity in the space
of local functionals, and that we redefine the action as described above
to make $G_1=(S,\Gamma_1)=0$. An anomaly that violates the master equation
$(\Gamma,\Gamma)=0$ in two-loop order is represented by a function $G_2$
for which
\[
 (\Gamma_1,\Gamma_1)+2(S,\Gamma_2)=G_2.
\]
But since $(S,\Gamma_1)=0$ and $(S,S)=0$, any such $G_2$ would satisfy
$(S,G_2)=0$. On the assumption of an empty cohomology, this means that it
can be expressed as $G_2=(S,F_2)$ with $F_2$ a local functional of ghost
number zero, so in this order the anomaly can be cancelled by subtracting
$F_2$ from the action.

This argument can be extended to all orders. Suppose we have cancelled
the anomalies in the master equation up to order $N-1$, so that
\[
 0=G_M
 =\sum_{L=0}^{M}(\Gamma_L,\Gamma_{M-L})
 =2(S,\Gamma_M)
 +\sum_{L=1}^{M-1}(\Gamma_L,\Gamma_{M-L}),
\]
for all $M<N$. The $N$th-order term in the antibracket
$(\Gamma,\Gamma)$ is likewise
\[
 G_N=2(S,\Gamma_N)
 +\sum_{M=1}^{N-1}(\Gamma_M,\Gamma_{N-M}),
\]
so, using the Jacobi identity \eqref{eq:15.9.21} (in which for three
bosonic operators all signs are $-$) and the above formula for
$(S,\Gamma_M)$, we find:
\[
 (S,G_N)
 =-2\sum_{M=1}^{N-1}\bigl((S,\Gamma_M),\Gamma_{N-M}\bigr)
 =\sum_{M=1}^{N-1}\sum_{L=1}^{M-1}
 \bigl((\Gamma_L,\Gamma_{M-L}),\Gamma_{N-M}\bigr).
\]
This can be written in a more symmetric way
\[
 (S,G_N)
 =-\sum_{M_1=1}^{N-2}\sum_{M_2=1}^{N-2}\sum_{M_3=1}^{N-2}
 \delta_{N,M_1+M_2+M_3}
 \bigl((\Gamma_{M_1},(\Gamma_{M_2},\Gamma_{M_3}))\bigr).
\]
Since the ranges of $M_1$, $M_2$, and $M_3$ are the same, we can write
the double antibracket in this sum as a sum over the $3!$ permutations
over these indices, which vanishes according to the Jacobi identity
\eqref{eq:15.9.21}, leading to the conclusion that $(S,G_N)=0$. If as
assumed the cohomology is empty then this implies that there is a local
functional $F_N$ for which $G_N=(S,F_N)$, so by subtracting $F_N$ from
the action the anomaly can be cancelled in order $N$, as was to be shown.

Using purely algebraic methods Barnich, Brandt, and
Henneaux~\hyperref[ch22-ref-22]{[22]} have succeeded in showing that for
Yang--Mills theories (in four spacetime dimensions) based on semisimple
gauge groups, the cohomology of the antibracket operation
$X\mapsto(S,X)$ (at ghost number unity on the space of local functionals)
consists entirely of a linear combination of terms of the form
\eqref{eq:22.6.20}, one for each simple subgroup of the gauge group, with
unknown coefficients.\footnote{It is not necessary to specify the
representation of the gauge algebra in which the trace in
Eq.~\eqref{eq:22.6.20} is to be calculated, because this trace is the same
up to a constant coefficient for all representations of a simple Lie
algebra,\hyperref[ch22-ref-19a]{[19a]} and the constant coefficient is not
determined anyway by this cohomology theorem.} This shows without any
reference to the matter content of the theory that in one-loop order,
where the anomaly $G_1$ automatically satisfies $(S,G_1)=0$, the anomaly
for a semi-simple gauge group must be a linear combination of terms of the
of the form \eqref{eq:22.6.20}, with only the constant coefficient for each
simple subgroup left to be determined by detailed calculations that take
into account the matter content of the theory. Further, we saw in
\hyperref[sec:22.4]{Section 22.4} that there are gauge groups for which the
trace in Eq.~\eqref{eq:22.6.20} vanishes automatically for any menu of
fermion fields. (These are the semisimple gauge groups with no $SU(n)$
factors with $n\geq3$.) In such cases the theorem of
Ref.~\hyperref[ch22-ref-22]{[22]} shows that the cohomology of the
antibracket operation $X\mapsto(S,X)$ at ghost number unity is zero. As we
have seen, this means that in such theories there is no anomaly in any
order of perturbation theory.

There is another sense in which the anomaly is related to the cohomology
of the antibracket operation~\hyperref[ch22-ref-23]{[23]}. In deriving the
Slavnov--Taylor identity \eqref{eq:16.4.6}, we assumed that the measure
$\prod_{n,x}d\chi^n(x)$ is invariant under the symmetry transformation in
question. In \hyperref[sec:15.9]{Section 15.9} the Zinn--Justin equation was
derived from the Slavnov--Taylor identity for the symmetry transformation
$\chi^n\rightarrow\chi^n+\theta\,\delta S/\delta\chi_n^*$, so the
derivation of the Zinn--Justin equation given in Section 15.9 breaks down
unless $\prod_{n,x}d\chi^n(x)$ is invariant under this transformation, or
in other words unless $\Delta S=0$, where $\Delta$ is the operator
\eqref{eq:15.9.34}. Where $\Delta S\neq0$, it still may be possible to save
the Zinn--Justin equation by adding local functionals to $S$ that violate
the classical master equation $(S,S)=0$ in such a way as to cancel the
effect of the non-invariance of the measure. It turns out that the
condition for this cancellation is nothing but the quantum master equation
\eqref{eq:15.9.35}. To construct an action $S$ that satisfies this
equation, we start with a zeroth-order action $S_0$ that satisfies the
classical master equation, $(S_0,S_0)=0$, and add quantum corrections. In
the case in which the cohomology of the operation $X\mapsto(S_0,X)$ (at
ghost number unity in the space of local functionals) is empty, the same
proof that was used above to show the absence in this case of anomalies to
all orders can be used to show that a local functional can be added to
$S_0$ so that the quantum master equation is satisfied to all orders.
